% ----------------------------------
%   - Ito's lemma (formula)
%       ? Some heuristics
%       ? Justification of existence the Ito integral
%   - Stochastic differential equations (SDEs
% ----------------------------------

\section{Stochastic calculus}

To model the dynamics of a financial asset's value,
an ordinary differential equation of the form
\begin{equation*}
    \frac{dS_t}{dt} = f(t, S_t)
\end{equation*}
is insufficient: it cannot capture the random fluctuations inherent in financial data.
Stochastic differential equations (SDEs) add a noise term to handle this,
\begin{equation}\label{eq:sde}
    dS_t = f(t, S_t)\, dt + \sigma(t, S_t)\, dW_t,
\end{equation}
where $W_t$ is a Brownian motion.
In integral form this reads
\begin{equation}\label{eq:int_sde}
    S_t = S_0 + \int^{t}_{0} f(s, S_s)\, ds + \int^{t}_{0} \sigma(s, S_s)\, dW_s,
\end{equation}
thus a key step is to give rigorous meaning to the stochastic integral $\int \sigma\, dW_s$.

The most well known process that has a well defined stochastic integral is the Wiener process or Brownian motion,
this stochastic integral is known as the Ito integral and will allow us to have the machinery to study SDEs and thus model the dynamics of financial assets.

\subsection{Ito's integral}

\begin{definition}\label{def:ito_integral_reqs}
    Let $\mathcal{V} = \mathcal{V}(S, T)$ the class of functions
    \begin{equation*}
        f(t, \omega): [0, \infty) \times \Omega \to \R,
    \end{equation*}
    such that
    \begin{itemize}
        \item $(t, \omega)$ is $\mathcal{B} \times \mathcal{F}$-measurable, where $\mathcal{B}$ is the Borel $\sigma$-algebra on $[0, \infty)$ and $\mathcal{F}$ is the $\sigma$-algebra on $\Omega$,
        \item $f(t, \omega)$ is $\mathcal{F}_t$ adapted.
        \item $\E \left[ \int^{T}_{S} f(t, \omega)^2 dt \right] < \infty$.
    \end{itemize}
\end{definition}

We will start by defining the Itô integral for \textit{elementary functions},
then this can be extended to functions in $\mathcal{V}$.

\begin{definition}
    A function $\phi$ is called \textit{elementary} if it has the form
    \begin{equation*}
        \phi(t, \omega) = \sum_{j} e_j(\omega) \1_{[t_j, t_{j + 1})}(t).
    \end{equation*}
\end{definition}

For elementary functions we define the Ito integral as
\begin{equation*}
    \int^{T}_{S} \phi(t, \omega)\, dB_t(\omega) = \sum_{j} e_j(\omega) \left[ B_{t_{j + 1}} - B_{t_{j}} \right] (\omega).
\end{equation*}

\begin{lemma}[Ito isometry]
    Let $\phi(t, \omega)$ be a bounded elementary function, then
    \begin{equation*}
        \E \left[ \left( \int^{T}_{S} \phi(t, \omega) dB_t (\omega) \right)^2 \right]
            = \E \left[ \int^{T}_{S} \phi(t, \omega)^2 dt \right].
    \end{equation*}
\end{lemma}

Using the Ito isometry and the definition of the Ito integral for elementary functions,
the Ito integral can be extended to the class of functions $\mathcal{V}$,
using that the functions in $\mathcal{V}$ can be expressed as the limit of sequences of simpler functions.
It is worth mentioning that the Itô isometry still holds for any function in $\mathcal{V}$.

The result of the integral is a stochastic process,
which is a continuous martingale.
This is similar to regular integration, where the result of integrating a function is a new function,
in this case the result of integrating a process "over" a stochastic process is a new stochastic process,
plus the resulting process has the nice property of being a martingale.
% Theorem? 6.2 Steele

The class $\mathcal{V}$ is rather limited and does not even contain all continuous functions, see \cite{steele_stochastic}.
Having to restrict to this class would be rather inconvenient when modeling using SDEs,
luckily the Itô integral can be extended to a broader class by relaxing the last condition of Definition~\ref{def:ito_integral_reqs}
\begin{equation*}
    P \left( \int^{T}_{S} f(t, \omega)^2 dt < \infty \right) = 1,
\end{equation*}
will call this class of functions $\mathcal{W}$.
Notice that the class $\mathcal{W}$ does contain all continuous functions.

\subsection{Ito's formula}

Having defined the Ito integral,
would be interesting to have a way to easily compute stochastic integrals without having to compute the limit of elementary functions,
that is the purpose of the Ito's lemma or Ito's formula.
The Ito's formula is a stochastic version of the chain rule for derivatives.

\begin{definition}[Ito process]
    Let $B_t$ be a Brownian motion on $(\Omega, \F, P)$.
    An Ito process is a stochastic process $X_t$ on $(\Omega, \F, P)$ of the form
    \begin{equation}\label{eq:ito_process}
        X_t = X_0 + \int^{t}_{0} \mu(s, X_s) ds + \int^{t}_{0} \sigma(s, X_s) dB_s,
    \end{equation}
    where $\mu$ and $\sigma$ are adapted, measurable processes, that satisfy the integrability conditions
    \begin{equation*}
        P \left( \int^{T}_{0} \left| \mu(t, \omega) \right| dt < \infty \right) = 1,
    \end{equation*}
    and
    \begin{equation*}
        P \left( \int^{T}_{0} \sigma(t, \omega)^2 dt < \infty \right) = 1.
    \end{equation*}
\end{definition}

The equation \eqref{eq:ito_process} can be also written in differential form as
\begin{equation*}
    dX_t = \mu\, dt + \sigma\, dB_t.
\end{equation*}

\begin{theorem}[1-dimensional Itô's formula]\label{thm:ito_formula}
    Let $X_t$ be an Itô process given by
    \begin{equation*}
        dX_t = \mu\, dt + \sigma\, dB_t.
    \end{equation*}
    Let $g(t, x) \in C^2([0, \infty) \times \R)$, then
    \begin{equation*}
        Y_t = g(t, X_t),
    \end{equation*}
    is an Ito process given by
    \begin{equation}\label{eq:ito_formula}
        dY_t = 
            \frac{\partial g}{\partial t}(t, X_t) dt
            + \frac{\partial g}{\partial x}(t, X_t) dX_t
            + \frac{1}{2} \frac{\partial^2 g}{\partial x^2}(t, X_t) (dX_t)^2,
    \end{equation}
    where $(dX_t)^2$ is computed according to the rules
    \begin{equation*}
        (dB_t)^2 = dt, \qquad
        (dt)^2 = 0, \qquad
        dB_t \cdot dt = 0.
    \end{equation*}
\end{theorem}

For the proof of the Ito's formula we refer to \cite{oksendal_sde}.
The formula can be extended to multiple dimensions.

\begin{definition}[$n$-dimensional Ito process]
    Let $B_t = (B_1(t), \dots, B_m(t))$ be a $m$-dimensional Brownian motion,
    let $\mu_{i}$ and $\sigma_{ij}$ be adapted, measurable processes such that
    \begin{equation*}
        P \left( \int^{T}_{0} \left| \mu_{i}(t, \omega) \right| dt < \infty \right) = 1,
        \qquad
        P \left( \int^{T}_{0} \sigma_{ij}(t, \omega)^2 dt < \infty \right) = 1,
    \end{equation*}
    then we define the $n$-dimensional Ito process $X_t = (X_1(t), \dots, X_n(t))$ as
    \begin{equation*}
        \left\{
        \begin{aligned}
            & dX_1(t) = \mu_1 dt + \sigma_{11} dB_1(t) + \dots + \sigma_{1m} dB_m(t), \\
            & \vdots \\
            & dX_n(t) = \mu_n dt + \sigma_{n1} dB_1(t) + \dots + \sigma_{nm} dB_m(t), \\
        \end{aligned}
        \right.
    \end{equation*}
    or in matrix notation
    \begin{equation*}
        dX_t = \mu\, dt + \sigma\, dB_t,
    \end{equation*}
    where $\mu = (\mu_1, \dots, \mu_n)$ and $\sigma = (\sigma_{ij})$ are an $n$-dimensional vector and an $n \times m$ matrix respectively.
\end{definition}

\begin{theorem}[General Itô's formula]\label{thm:ndim_ito_formula}
    Let $X_t$ be an $n$-dimensional Itô process given by
    \begin{equation*}
        dX_t = \mu\, dt + \sigma\, dB_t.
    \end{equation*}
    Let $g(t, x) = (g_1(t, x), \dots, g_p(t, x)) \in C^2([0, \infty) \times \R^n)$, then
    \begin{equation*}
        Y_t = g(t, X_t),
    \end{equation*}
    is a $p$-dimensional Ito process whose components are given by
    \begin{equation}\label{eq:ndim_ito_formula}
        dY_k(t) = 
            \frac{\partial g_k}{\partial t}(t, X_t) dt
            + \sum_{i} \frac{\partial g_k}{\partial x_i}(t, X_t) dX_i(t)
            + \frac{1}{2} \sum_{ij} \frac{\partial^2 g_k}{\partial x_i \partial x_j}(t, X_t) dX_i(t) dX_j(t),
    \end{equation}
    where $dX_i(t) dX_j(t)$ is computed according to the rules
    \begin{equation*}
        (dB_i(t))^2 = dt, \qquad
        dB_i(t)\, dB_j(t) = 0 \quad (i \neq j).
    \end{equation*}
\end{theorem}

\subsection{Stochastic differential equations}

A \textbf{stochastic differential equation} (SDE) is an equation of the form
\begin{equation}\label{eq:sde_general}
    dX_t = b(t, X_t)\, dt + \sigma(t, X_t)\, dB_t, \qquad X_0 = x_0,
\end{equation}
where $b: [0,T] \times \R^n \to \R^n$ is the \textit{drift} and
$\sigma: [0,T] \times \R^n \to \R^{n \times m}$ is the \textit{diffusion} coefficient.
A solution is an Itô process $X_t$ satisfying \eqref{eq:sde_general} almost surely.

The following theorem, whose proof may be found in \cite[Thm.~5.2.1]{oksendal_sde},
gives sufficient conditions for a solution to exist and be unique.

\begin{theorem}[Existence and uniqueness]\label{thm:sde_existence}
    Let $T > 0$.
    Suppose $b$ and $\sigma$ satisfy, for all $t \in [0,T]$ and $x, y \in \R^n$,
    \begin{enumerate}
        \item \textup{(Lipschitz)} $|b(t,x) - b(t,y)| + |\sigma(t,x) - \sigma(t,y)| \leq C\,|x - y|$,
        \item \textup{(Linear growth)} $|b(t,x)|^2 + |\sigma(t,x)|^2 \leq C^2(1 + |x|^2)$,
    \end{enumerate}
    for some constant $C > 0$, and let $x_0 \in \R^n$.
    Then the SDE \eqref{eq:sde_general} has a unique (up to indistinguishability)
    continuous adapted solution $\{X_t\}_{0 \leq t \leq T}$.
\end{theorem}

We close with two fundamental examples, both following \cite[Ch.~5]{oksendal_sde}.

\subsubsection{Example: Geometric Brownian motion}

Let $\mu, \sigma \in \R$ with $\sigma \neq 0$ and $S_0 > 0$.
Consider the SDE
\begin{equation}\label{eq:gbm_sde}
    dS_t = \mu S_t\, dt + \sigma S_t\, dB_t.
\end{equation}
This is called \textbf{geometric Brownian motion} and is the model used in Section~\ref{subsec:black_scholes_model}
for the dynamics of a stock price.

Applying Theorem~\ref{thm:ito_formula} to $g(t, x) = \ln x$,
with $g_t = 0$, $g_x = x^{-1}$, $g_{xx} = -x^{-2}$,
\begin{equation*}
    d(\ln S_t)
        = \frac{1}{S_t}\, dS_t - \frac{1}{2 S_t^2}(dS_t)^2
        = \mu\, dt + \sigma\, dB_t - \frac{\sigma^2}{2}\, dt
        = \left(\mu - \frac{\sigma^2}{2}\right) dt + \sigma\, dB_t.
\end{equation*}
Integrating both sides from $0$ to $t$ gives $\ln(S_t / S_0) = (\mu - \sigma^2/2)\,t + \sigma B_t$,
so the unique solution is
\begin{equation}\label{eq:gbm_solution}
    S_t = S_0 \exp\!\left[\left(\mu - \frac{\sigma^2}{2}\right) t + \sigma B_t\right].
\end{equation}

\subsubsection{Example: Ornstein-Uhlenbeck process}

Let $\theta, \sigma > 0$ and $\mu \in \R$.
Consider the SDE
\begin{equation}\label{eq:ou_sde}
    dX_t = \theta(\mu - X_t)\, dt + \sigma\, dB_t, \qquad X_0 = x_0.
\end{equation}
This is the \textbf{Ornstein-Uhlenbeck (OU) process},
a mean-reverting model: the drift $\theta(\mu - X_t)$ pulls the process toward the long-run mean $\mu$ at rate $\theta$.
It appears, for instance, in the Vasicek interest rate model.

To solve, let $Y_t = (X_t - \mu)e^{\theta t}$.
Applying Theorem~\ref{thm:ito_formula} with $g(t,x) = (x - \mu)e^{\theta t}$,
\begin{equation*}
    dY_t
        = \theta(X_t - \mu)e^{\theta t}\, dt + e^{\theta t}\, dX_t
        = e^{\theta t}\bigl[\theta(\mu - X_t) + \theta(X_t - \mu)\bigr] dt + \sigma e^{\theta t}\, dB_t
        = \sigma e^{\theta t}\, dB_t.
\end{equation*}
Integrating and reverting the substitution,
\begin{equation}\label{eq:ou_solution}
    X_t = \mu + (x_0 - \mu)\,e^{-\theta t} + \sigma \int_0^t e^{-\theta(t - s)}\, dB_s.
\end{equation}
The OU process is Gaussian with $\E[X_t] = \mu + (x_0 - \mu)e^{-\theta t} \to \mu$ as $t \to \infty$.




